\begin{ex}
Giảm 24 đi 3 lần ta được số nào?
\shortans{8}
\loigiai{
$24 : 3 = 8$.
}
\end{ex}

\begin{ex}
Giảm 36 đi 4 lần ta được số nào?
\shortans{9}
\loigiai{
$36 : 4 = 9$.
}
\end{ex}

\begin{ex}
Giảm 45 đi 5 lần ta được số nào?
\shortans{9}
\loigiai{
$45 : 5 = 9$.
}
\end{ex}

\begin{ex}
Giảm 48 đi 6 lần ta được số nào?
\shortans{8}
\loigiai{
$48 : 6 = 8$.
}
\end{ex}

\begin{ex}
Giảm 56 đi 7 lần ta được số nào?
\shortans{8}
\loigiai{
$56 : 7 = 8$.
}
\end{ex}

\begin{ex}
Giảm 64 đi 8 lần ta được số nào?
\shortans{8}
\loigiai{
$64 : 8 = 8$.
}
\end{ex}

\begin{ex}
Giảm 72 đi 9 lần ta được số nào?
\shortans{8}
\loigiai{
$72 : 9 = 8$.
}
\end{ex}

\begin{ex}
Giảm 80 đi 4 lần ta được số nào?
\shortans{20}
\loigiai{
$80 : 4 = 20$.
}
\end{ex}

\begin{ex}
Giảm 90 đi 3 lần ta được số nào?
\shortans{30}
\loigiai{
$90 : 3 = 30$.
}
\end{ex}

\begin{ex}
Một sợi dây dài 28 cm. Sau khi cắt đi thì độ dài sợi dây giảm đi 4 lần. Hỏi sợi dây còn lại dài bao nhiêu xăng-ti-mét?
\shortans{7}
\loigiai{
$28 : 4 = 7$ cm.
}
\end{ex}

\begin{ex}
Bác An thu hoạch được 60 kg cam. Bác Ba thu hoạch được số cam giảm đi 3 lần so với bác An. Hỏi bác Ba thu hoạch được bao nhiêu ki-lô-gam cam?
\shortans{20}
\loigiai{
$60 : 3 = 20$ kg.
}
\end{ex}

\begin{ex}
Năm nay mẹ 36 tuổi, tuổi con bằng tuổi mẹ giảm đi 4 lần. Hỏi năm nay con bao nhiêu tuổi?
\shortans{9}
\loigiai{
$36 : 4 = 9$ tuổi.
}
\end{ex}

\begin{ex}
Bố 42 tuổi, con có tuổi bằng tuổi bố giảm đi 6 lần. Hỏi con bao nhiêu tuổi?
\shortans{7}
\loigiai{
$42 : 6 = 7$ tuổi.
}
\end{ex}

\begin{ex}
Một cuộn vải dài 70 m. Sau khi bán, số vải giảm đi 5 lần. Hỏi cuộn vải còn lại bao nhiêu mét?
\shortans{14}
\loigiai{
$70 : 5 = 14$ m.
}
\end{ex}

\begin{ex}
Giảm 60 đi 5 lần rồi cộng thêm 18 ta được số nào?
\shortans{30}
\loigiai{
$60 : 5 + 18 = 12 + 18 = 30$.
}
\end{ex}

\begin{ex}
Giảm 84 đi 4 lần rồi bớt đi 11 ta được số nào?
\shortans{10}
\loigiai{
$84 : 4 - 11 = 21 - 11 = 10$.
}
\end{ex}

\begin{ex}
Một số giảm đi 6 lần thì được 8. Số đó là số nào?
\shortans{48}
\loigiai{
$8 \times 6 = 48$.
}
\end{ex}

\begin{ex}
Một số giảm đi 7 lần thì được 9. Số đó là số nào?
\shortans{63}
\loigiai{
$9 \times 7 = 63$.
}
\end{ex}

\begin{ex}
Thùng thứ nhất có 48 lít dầu. Thùng thứ hai có số dầu giảm đi 4 lần so với thùng thứ nhất. Hỏi thùng thứ hai có bao nhiêu lít dầu?
\shortans{12}
\loigiai{
$48 : 4 = 12$ lít.
}
\end{ex}

\begin{ex}
Một rổ có 54 quả lê. Số quả táo ít hơn số quả lê giảm đi 6 lần. Hỏi có bao nhiêu quả táo?
\shortans{9}
\loigiai{
$54 : 6 = 9$ quả táo.
}
\end{ex}

\begin{ex}
Một số giảm đi 3 lần thì bằng 15. Tìm số đó.
\shortans{45}
\loigiai{
$15 \times 3 = 45$.
}
\end{ex}

\begin{ex}
Một người có 96 con gà. Sau khi bán thì số gà giảm đi 3 lần. Hỏi người đó còn lại bao nhiêu con gà?
\shortans{32}
\loigiai{
$96 : 3 = 32$ con gà.
}
\end{ex}

\begin{ex}
Giảm 50 đi 5 lần được bao nhiêu?
\shortans{10}
\loigiai{
$50 : 5 = 10$.
}
\end{ex}

\begin{ex}
Giảm 100 đi 4 lần được bao nhiêu?
\shortans{25}
\loigiai{
$100 : 4 = 25$.
}
\end{ex}

\begin{ex}
Đoạn thẳng AB dài 32 cm. Đoạn CD có độ dài bằng độ dài AB giảm đi 4 lần. CD dài bao nhiêu xăng-ti-mét?
\shortans{8}
\loigiai{
$32 : 4 = 8$ cm.
}
\end{ex}
